\documentclass[11pt]{article}

\usepackage[T1]{fontenc}
\usepackage[utf8]{inputenc}
\usepackage{lmodern}
\usepackage{amsmath,amssymb,amsthm,mathtools}
\usepackage{booktabs}
\usepackage{geometry}
\usepackage{microtype}
\usepackage[hidelinks]{hyperref}

\geometry{margin=1in}
\setlength{\emergencystretch}{2em}
\hypersetup{
  pdftitle={A Certified Zero Count in a Finite Fibonacci-Zeta Rectangle},
  pdfauthor={Xiaojian Zeng}
}

\newtheorem{theorem}{Theorem}
\newtheorem{lemma}[theorem]{Lemma}
\newtheorem{proposition}[theorem]{Proposition}
\theoremstyle{remark}
\newtheorem{remark}[theorem]{Remark}

\newcommand{\C}{\mathbb{C}}
\newcommand{\ii}{\mathrm{i}}
\newcommand{\ZF}{Z_F}
\newcommand{\real}{\operatorname{Re}}
\newcommand{\imag}{\operatorname{Im}}

\title{A Certified Zero Count in a Finite Fibonacci-Zeta Rectangle}
\author{Xiaojian Zeng\thanks{PBC School of Finance, Tsinghua University,
100190, China. Email: \href{mailto:xiaojianzeng@foxmail.com}{xiaojianzeng@foxmail.com}.}}
\date{8 September 2026}

\begin{document}
\maketitle

\begin{abstract}
Let $F_1=F_2=1$, $F_{n+2}=F_{n+1}+F_n$, and
$\ZF(s)=\sum_{n\geq 1}F_n^{-s}$ for $\real s>0$.  We prove, with zeros
counted with analytic multiplicity, that
\[
 \#\{\rho:\ZF(\rho)=0,\ \real\rho\geq0.73,\
                 0<\imag\rho\leq8596696\}=1.
\]
Moreover, the closed disk of radius $10^{-6}$ centered at
$0.730178633163+8596695.206382\ii$ contains exactly one distinct zero,
and that zero is simple.  The proof combines a published phase-locking
necessity with an exhaustive outward-rounded ordinate cover, uniform
Fibonacci-series tail bounds, a closed-box nonvanishing cover, and a strict
Rouch\'e comparison.  The argument treats every boundary included in the
count; it is not an extrapolation from a targeted zero search.
\end{abstract}

\section{Statement and context}

The Dirichlet series $\ZF$ converges absolutely in $\real s>0$.  Mantovanelli
determined the right edge of the closure of its zero real parts
\cite{Mantovanelli2026}.  In particular, there is a number
\[
 \sigma_F=0.743163398726901648\ldots
\]
such that $\ZF$ has no zero in $\real s\geq\sigma_F$.  The same work gives
quantitative necessary phase-locking conditions below this edge and reports
a zero numerically near
$0.730178633163+8596695.206382\ii$.  That reported search is explicitly
targeted rather than exhaustive \cite[Section~7]{Mantovanelli2026}; hence it
does not determine the number of zeros in the finite region considered here.

For $T>0$, set
\[
 N_{0.73}(T)=\#\{\rho:\ZF(\rho)=0,\ \real\rho\geq0.73,\
                                  0<\imag\rho\leq T\},
\]
where analytic multiplicity is used.  Let
\[
 c=0.730178633163+8596695.206382\ii,
 \qquad r=10^{-6},
 \qquad D=\{s\in\C:|s-c|\leq r\}.
\]

\begin{theorem}\label{thm:main}
One has $N_{0.73}(8596696)=1$.  The closed disk $D$ contains exactly one
distinct zero of $\ZF$, and that zero is simple.
\end{theorem}

This is a finite census only.  It makes no assertion about zeros above the
stated height, an exact first-hitting function, or the shrinking-target
heuristic in \cite{Mantovanelli2026}.  The proof is computer-assisted, but
all numerical predicates below are strict interval predicates and the
complete finite covers are serialized with the accompanying programs.

\section{The phase-lock input}

We first record exactly what is imported from \cite{Mantovanelli2026}.  For
$\sigma\geq\tfrac12$, write
\begin{align*}
 A&=2^{-\sigma},&B&=3^{-\sigma},&C&=8^{-\sigma},&D_0&=144^{-\sigma},\\
 L(\sigma)&=2-A-B-C+D_0,
 &R(\sigma)&=\sum_{\substack{n\geq5\\n\neq6,12}}F_n^{-\sigma},
 &H(\sigma)&=R(\sigma)-L(\sigma),\\
 K(\sigma)&=72^{-\sigma}+18^{-\sigma}+4\,48^{-\sigma}.
\end{align*}
Theorem~1.1 and Lemma~2.1 of that paper state that $\sigma_F\in(0.7,0.75)$,
$H$ is decreasing with $H(\sigma_F)=0$, and $\ZF(s)\neq0$ for
$\real s\geq\sigma_F$.  Its Proposition~7.1 states the following.  If
$\rho=\sigma+\ii t$ is a zero and $\tfrac12\leq\sigma<\sigma_F$, then
\begin{equation}\label{eq:17}
 A|2^{-\ii t}+1|^2+B|3^{-\ii t}+1|^2
       \leq \frac{2H(\sigma)}{1-K(\sigma)}.
\end{equation}
If
\[
 E(\rho)=\sum_{\substack{n\geq5\\n\neq6,12}}F_n^{-\sigma}F_n^{-\ii t},
 \qquad \omega=\frac{E(\rho)}{|E(\rho)|},
\]
then
\begin{align}
 \sum_{\substack{n\geq5\\n\neq6,12}}
 F_n^{-\sigma}|F_n^{-\ii t}-\omega|^2&\leq2H(\sigma),
                                                        \label{eq:18}\\
 |\omega+1|^2&\leq\frac{2H(\sigma)}{L(\sigma)+H(\sigma)},
                                                        \label{eq:19}\displaybreak[0]\\
 |F_n^{-\ii t}+1|&\leq
 \sqrt{2H(\sigma)F_n^\sigma}
 +\sqrt{\frac{2H(\sigma)}{L(\sigma)+H(\sigma)}}
                                                        \label{eq:20}
\end{align}
for every $n\geq5$ with $n\neq6,12$.  These are necessary conditions; no
existence conclusion is inferred from them.

\section{An exhaustive ordinate cover}

\begin{lemma}[Certified phase cover]\label{lem:phase-cover}
If $\ZF(\sigma+\ii t)=0$ with
\[
 0.73\leq\sigma<\sigma_F,
 \qquad 0<t\leq8596696,
\]
then $t$ belongs to one of the eight closed intervals in
Table~\ref{tab:phase-intervals}.
\end{lemma}

\begin{table}[ht]
\centering
\caption{The certified necessary cover of possible ordinates.}
\label{tab:phase-intervals}
\begin{tabular}{r@{\qquad}l}
\toprule
$j$ & $I_j$\\
\midrule
1 & $[893451.166593898,\ 893451.296590004]$\\
2 & $[1045810.765660011,\ 1045810.914566694]$\\
3 & $[1292959.408434021,\ 1292959.524502041]$\\
4 & $[1431405.026953860,\ 1431405.066784306]$\\
5 & $[2302570.625261012,\ 2302570.893936211]$\\
6 & $[2396445.266260887,\ 2396445.372661887]$\\
7 & $[2688164.836889199,\ 2688164.901257748]$\\
8 & $[8596695.022433829,\ 8596695.288582920]$\\
\bottomrule
\end{tabular}
\end{table}

\begin{proof}
We give both the mathematical reduction and the finite certification rule.
Throughout $0.73\leq\sigma<\sigma_F<0.75$, monotonicity and directed ball
evaluation give the strict bounds
\begin{equation}\label{eq:phase-constants}
\begin{gathered}
 H(\sigma)<0.045602822,\qquad K(\sigma)<0.402314890,\\
 2^{-\sigma}>0.594603557,\qquad
 3^{-\sigma}>0.438691337,\qquad
 R(\sigma)>0.75746608.
\end{gathered}
\end{equation}
For the last inequality it suffices to use the terms through $F_{90}$ at
$\sigma=3/4$.  All terminating decimals in \eqref{eq:phase-constants} are
exact rationals on the safe side of the indicated inequality.

If $|e^{-\ii x}+1|\leq q<2$, then the distance of $x$ from an odd multiple
of $\pi$, modulo $2\pi$, is at most $2\arcsin(q/2)$.  Applying this observation
to \eqref{eq:17} and \eqref{eq:20}, using
\eqref{eq:phase-constants}, gives odd-$\pi$ windows for the bases shown in
Table~\ref{tab:angular-bounds}.  The displayed numbers are the exact rational
upper bounds used for their angular radii.  The bound from \eqref{eq:20}
first becomes vacuous at $F_{13}=233$ and is increasing thereafter, so the
listed free indices are all the nonvacuous ones.

\begin{table}[ht]
\centering
\caption{Outward upper bounds $\delta_b$ in
$\operatorname{dist}(t\log b,(2\mathbb Z+1)\pi)\leq\delta_b$.}
\label{tab:angular-bounds}
\small
\begin{tabular}{r@{\quad}l r@{\quad}l}
\toprule
$b$&$\delta_b$&$b$&$\delta_b$\\
\midrule
2  &$0.512174793486796593149426$&3  &$0.598687998441899467344934$\\
5  &$0.932674913662750792270648$&13 &$1.209611441366676873107514$\\
21 &$1.405838206058152966984497$&34 &$1.666472314236621336794858$\\
55 &$2.039987021672562880879150$&89 &$2.810571796660721292073878$\\
\bottomrule
\end{tabular}
\end{table}

Here is the endpoint certification.  Given outward rational enclosures
$\pi^-<\pi<\pi^+$ and $\ell_b^-<\log b<\ell_b^+$, the window belonging to
an odd positive integer $m$ is contained in
\begin{equation}\label{eq:periodic-window}
 \left[
   \frac{m\pi^- -\delta_b}{\ell_b^+},
   \frac{m\pi^+ +\delta_b}{\ell_b^-}
 \right].
\end{equation}
The program checks 24-decimal enclosures for $\pi$ and every logarithm by
192-bit ball arithmetic, rounds both endpoints of \eqref{eq:periodic-window}
outward to the $10^{-9}$ grid, and intersects the eight sorted periodic
covers using integer arithmetic.  After the bases $2,3$ there are $425784$
intervals; after successively imposing the bases $5,13,21,34,55,89$, the
counts are $194823,107363,62843,41004,31269,31425$.  At this stage the total
width is $4193.521920257$.

It remains to describe the certified joint pruning, which is essential for
obtaining the short final list.  Put
\[
 w_n=F_n^{-3/4},\qquad
 S(t)=\sum_{\substack{5\leq n\leq30\\n\neq6,12}}w_nF_n^{-\ii t},
 \qquad W=\sum_{\substack{5\leq n\leq30\\n\neq6,12}}w_n.
\]
Because $w_n\leq F_n^{-\sigma}$, equation \eqref{eq:18} implies
\begin{equation}\label{eq:joint-necessary}
 W-\real\bigl(\overline\omega S(t)\bigr)\leq H(\sigma).
\end{equation}
Equation \eqref{eq:19} and \eqref{eq:phase-constants} allow us to write
$\omega=-e^{\ii\alpha}$ with
$|\alpha|\leq0.348764756$.  Consequently both
\begin{equation}\label{eq:joint-tests}
 W-|S(t)|\leq0.045602822,
 \qquad
 W-\max_{|\alpha|\leq0.348764756}
 \real\bigl(-e^{-\ii\alpha}S(t)\bigr)\leq0.045602822
\end{equation}
are necessary.

Each current ordinate interval is bisected until its width is at most
$0.005$.  On each box, the sines and cosines in $S(t)$ are evaluated as
balls.  The $\alpha$-range is covered by 24 closed subintervals, and an
interval upper bound is taken for every projection in the second part of
\eqref{eq:joint-tests}.  A box is discarded only if a strict interval
comparison disproves a necessary inequality in \eqref{eq:joint-tests}; every
inconclusive terminal box is retained in full.  There are $59741$ interval
tests: $42763$ boxes are excluded by the first test and $2476$ additional
boxes by the restricted-$\alpha$ test.  Merging the retained closed boxes
gives exactly the eight intervals in Table~\ref{tab:phase-intervals}, with
total width $1.140395094$.  Since every rounding is outward and every
discarded box violates a necessary condition, this is a cover rather than a
sample.
\end{proof}

\section{Analytic bounds and the closed-box certificate}

Let
\[
 P_N(s)=\sum_{n=1}^N F_n^{-s},\qquad N=140,
 \qquad \varphi=\frac{1+\sqrt5}{2}.
\]

\begin{lemma}[Uniform series bounds]\label{lem:analytic-bounds}
For $\real s\geq0.73$,
\begin{equation}\label{eq:uniform-bounds}
 |\ZF(s)-P_N(s)|<10^{-20},\qquad
 |\ZF'(s)-P_N'(s)|<10^{-18},
\end{equation}
and
\begin{equation}\label{eq:derivative-bounds}
 |\ZF'(s)|<3.79,\qquad |\ZF''(s)|<16.
\end{equation}
\end{lemma}

\begin{proof}
The elementary inequalities
\[
 \varphi^{n-2}\leq F_n\leq\varphi^{n-1}\qquad(n\geq2)
\]
follow by induction from $\varphi^2=\varphi+1$.  With
$q=\varphi^{-0.73}$ and $k_0=N-1=139$, they give
\begin{align*}
 \sum_{n>N}F_n^{-0.73}
 &\leq\frac{q^{k_0}}{1-q}
   <2.101\times10^{-21}<10^{-20},\\
 \sum_{n>N}(\log F_n)F_n^{-0.73}
 &\leq(\log\varphi)\frac{q^{k_0}
       ((k_0+1)-k_0q)}{(1-q)^2}
   <1.440\times10^{-19}<10^{-18}.
\end{align*}
This proves \eqref{eq:uniform-bounds}.  Directed evaluation of the first
$140$ positive terms at $\sigma=0.73$, plus the second displayed tail,
gives the strict upper bound $3.789075$ for $|\ZF'|$.  Finally,
\[
 |\ZF''(s)|
 \leq(\log\varphi)^2\sum_{k\geq1}(k+1)^2q^k
 = (\log\varphi)^2\frac{q(4-3q+q^2)}{(1-q)^3}
 <14.947983<16.
\]
All four final comparisons were evaluated with 256-bit balls.
\end{proof}

For a closed axis-parallel box
\[
 B=[\sigma_-,\sigma_+] + \ii[t_-,t_+],
\]
let $m(B)$ be its center and put
\[
 r_1(B)=\frac{(\sigma_+-\sigma_-)+(t_+-t_-)}2.
\]
Every $s\in B$ satisfies $|s-m(B)|\leq r_1(B)$.  Lemma
\ref{lem:analytic-bounds} therefore yields the following checkable
nonvanishing criterion:
\begin{equation}\label{eq:box-test}
 |P_N(m(B))|>10^{-20}+3.79r_1(B)
 \quad\Longrightarrow\quad
 \ZF(s)\neq0\quad(s\in B).
\end{equation}
Indeed, integrate $\ZF'$ along the segment from $m(B)$ to $s$ and use the
value-tail bound at the center.

\begin{lemma}[Closed-box cover]\label{lem:box-cover}
Every point in
\[
 \bigcup_{j=1}^8\bigl([0.73,0.75]+\ii I_j\bigr)
\]
either belongs to $D$ or lies in a closed box on which the strict test
\eqref{eq:box-test} holds.
\end{lemma}

\begin{proof}
Start with the eight closed rectangles in the statement.  For a current box,
first test whether all four corners lie in $D$; this is an exact rational
comparison, and convexity then places the whole box in $D$.  Otherwise,
evaluate $P_N$ at its rational center with 256-bit complex ball arithmetic.
If \eqref{eq:box-test} is certified strictly, retain the box with the label
``nonzero.''  If neither test succeeds, bisect the longer side at its exact
rational midpoint and repeat.  Closed children cover their parent, including
the common boundary.

The recursion visits $246$ boxes and terminates with the complete leaf counts
in Table~\ref{tab:leaf-counts}.  Thus $122$ leaves satisfy
\eqref{eq:box-test}, while the remaining five lie wholly in $D$.  The exact
rational endpoints and the strict lower modulus, upper error, and gap balls
for every leaf are recorded in the accompanying table.  This serialized
bisection tree also checks that no region between leaves is omitted.

\begin{table}[ht]
\centering
\caption{Leaves in the certified box cover.}
\label{tab:leaf-counts}
\begin{tabular}{rrrr}
\toprule
$j$ & nonzero leaves & leaves in $D$ & maximum depth\\
\midrule
1&1&0&0\\
2&1&0&0\\
3&1&0&0\\
4&1&0&0\\
5&4&0&3\\
6&1&0&0\\
7&1&0&0\\
8&112&5&36\\
\midrule
Total&122&5&36\\
\bottomrule
\end{tabular}
\end{table}
\end{proof}

\section{The disk certificate and completion of the count}

\begin{lemma}[Rouch\'e certificate]\label{lem:rouche}
The disk $D$ contains exactly one zero of $\ZF$, counted with multiplicity,
and $\ZF$ has no zero on $\partial D$.
\end{lemma}

\begin{proof}
Define the linear function
\[
 \mathcal L(s)=P_N(c)+P_N'(c)(s-c).
\]
The 256-bit ball evaluations at the exact rational center $c$ certify
\begin{align}
 |P_N(c)|&<1.500649\times10^{-6},&
 |P_N'(c)|&>3.272560,\label{eq:center-bounds}\\
 \frac{|P_N(c)|}{|P_N'(c)|}
   &<4.58555\times10^{-7}<r.&\label{eq:linear-zero}
\end{align}
Thus the unique zero of $\mathcal L$ lies in the interior of $D$.  On
$|s-c|=r$, the reverse triangle inequality gives the certified lower bound
\begin{equation}\label{eq:linear-boundary}
 |\mathcal L(s)|
 \geq |P_N'(c)|r-|P_N(c)|
 >1.7719116\times10^{-6}.
\end{equation}

Apply Taylor's formula to the full function $\ZF$ at $c$.  Equations
\eqref{eq:uniform-bounds}--\eqref{eq:derivative-bounds} give, on the same
circle,
\begin{align}
 |\ZF(s)-\mathcal L(s)|
 &\leq10^{-20}+10^{-18}r+\frac{16}{2}r^2\\
 &<8.000000011\times10^{-12}.
                                                        \label{eq:rouche-error}
\end{align}
The upper bound in \eqref{eq:rouche-error} is strictly smaller than the lower
bound in \eqref{eq:linear-boundary}.  Rouch\'e's theorem therefore gives the
same number of zeros, counted with multiplicity, for $\ZF$ and $\mathcal L$
inside the boundary.  That number is one.  The strict comparison also
excludes a boundary zero.
\end{proof}

\begin{proof}[Proof of Theorem~\ref{thm:main}]
Mantovanelli's zero-free theorem excludes $\real s\geq\sigma_F$.  By
Lemma~\ref{lem:phase-cover}, every remaining zero counted by
$N_{0.73}(8596696)$ has its ordinate in one of the $I_j$.  Lemma
\ref{lem:box-cover} excludes every point of the corresponding closed
rectangles outside $D$, while Lemma~\ref{lem:rouche} counts one zero in $D$.

Finally,
\[
 \real c-r>0.73,\qquad 0<\imag c-r,\qquad
 \imag c+r<8596696,
\]
so $D$ lies strictly inside the region being counted.  The phase and box
covers are closed, hence include $\real s=0.73$, all retained ordinate
endpoints, and any possible point on the top boundary.  It follows that the
full multiplicity-weighted count is one.  Total multiplicity one also means
that the disk contains one distinct zero and that this zero is simple.
\end{proof}

\section{Reproducibility and assistance disclosure}

The accompanying file \texttt{phase\_cover.py} certifies Lemma
\ref{lem:phase-cover} and writes the eight intervals and a machine-readable
report.  The file \texttt{zero\_count\_certificate.py} certifies Lemmas
\ref{lem:analytic-bounds}--\ref{lem:rouche} and writes the complete leaf
table and report.  From the project root, the deterministic commands are
\begin{verbatim}
python3 evidence/phase_cover.py
python3 evidence/zero_count_certificate.py
\end{verbatim}
The recorded environment used Python~3.10.12, python-flint~0.9.0, and
FLINT~3.6.0.  Transcendental and complex enclosures use Arb midpoint-radius
ball arithmetic \cite{Johansson2017Arb}; rational grid endpoints and box
bisections are exact.  Every theorem-bearing branch is fail-closed: a box is
removed only after a strict interval comparison, and otherwise it is retained
or subdivided.

For another explicit example of separating a finite reduction, exhaustive
coverage, boundary handling, and exact reconstruction of local checks, see
the certificate architecture in \cite{ZengEtAl2026FiniteCertificates}.  The
mathematical objects there are different; the relevant methodological point
is that coverage rather than a targeted search carries the universal finite
conclusion.

OpenAI Codex assisted with literature search, verification-program drafting
and auditing, LaTeX organization, and initial manuscript preparation.
Language-model output is not used as mathematical evidence.  The mathematical
claims above rest on the stated reductions, strict ball enclosures, and
serialized exhaustive covers.  Responsibility for any disseminated version
remains with the author.

\bibliographystyle{plain}
\bibliography{references}

\end{document}
